\chapter{最適化への応用}
    \section{狭義凸関数の線形制約下での最小化}
        \begin{shadebox}
            $m,n\in\naturalNumbers$とする。
            $f: \realNumbers^n \to \realNumbers$を狭義凸関数とする。
            $\bm{c}_i \in \realNumbers^n, b_i \in \realNumbers\;(i=1,\dotsc,m)$とし、$g_i: \bm{x} \in \realNumbers^n \mapsto \bm{c}_i^\top\bm{x} + b \in \realNumbers$とする。
            $\bm{\lambda} \in \realNumbers^m$とし、Lagrange関数$F$を$F(\bm{x},\bm{\lambda}) \coloneqq f(\bm{x}) + \sum_{i=1}^m \lambda_i g_i(\bm{x})$とする。
            $\bm{x}_0, \bm{\lambda}_0$を$F$の停留点、すなわち$\nabla_\bm{x}F(\bm{x}_0, \bm{\lambda}_0) = \bm{0}$とすると$f$は$g_i(\bm{x}) = 0\;(i=1,\dotsc,m)$なる拘束条件の下、$\bm{x}_0$で最小値をとる。
        \end{shadebox}
        \begin{proof}
            \quad\par
            $\bm{h} \in \realNumbers^n \neq \bm{0}$とし、$g_i(\bm{x} + \bm{h}) = 0\;(i=1,\dotsc,m)$を満たすとする。
            このとき$f(\bm{x}_0 + \bm{h}) > f(\bm{x}_0)$であることを示す。
            \[ f(\bm{x}_0 + \bm{h}) = f(\bm{x}_0) + \nabla f(\bm{x}_0)^\top\bm{h} + \bm{h}^\top H(f,\bm{x}_0)\bm{h} + o(\|\bm{h}\|_2^3) \]
            ここに、$H(f,\bm{x}_0)$は$f$のHesse行列の$\bm{x}_0$に於ける値である。
            $\nabla_\bm{x}F(\bm{x}_0, \bm{\lambda}_0) = \bm{0}$より$\nabla f(\bm{x}_0) = -\sum_{i=1}^m \bm{\lambda}_0[i] \nabla g_i(\bm{x}_0)$である。
            $g_i(\bm{x} + \bm{h}) = 0\;(i=1,\dotsc,m)$より$\bm{c}_i^\top\bm{h} = \bm{0}$であるから
            \[ \nabla f(\bm{x}_0)^\top\bm{h} = -\sum_{i=1}^m \bm{\lambda}_0[i] \nabla g_i(\bm{x}_0)\bm{h} = -\sum_{i=1}^m \bm{\lambda}_0[i] \bm{c}_i^\top\bm{h} = \bm{0} \]
            である。よって
            \[ f(\bm{x}_0 + \bm{h}) = f(\bm{x}_0) + \bm{h}^\top H(f,\bm{x}_0)\bm{h} + o(\|\bm{h}\|_2^3) \]
            $f$は狭義凸であるので$H(f,\bm{x}_0)$は正定であり、$O(\|\bm{h}\|_2^2)$であるから、十分小さい$\bm{h}$に対して上式の右辺は$f(\bm{x}_0)$を厳密に上回る。
        \end{proof}